% Source: https://artofproblemsolving.com/community/c3h2452011_ash_ketchum_problems
% Stable identity: aops-topic-2452011-q22
The lid of Ash's box is a regular hexagon with side length $69$. For each vertex $V$ of the hexagon, the arc of the circle with centre $V$ and radius $69$ that lies inside the hexagon is drawn; it runs from one neighbouring vertex of $V$ through the centre of the hexagon to the other neighbouring vertex. These six arcs bound six congruent leaf-shaped regions (Ash's tea leaves), each joining the centre of the hexagon to a vertex, as in the figure.

\begin{asy}
usepackage("amsmath");
size(4cm);
defaultpen(fontsize(10pt));
real r=1;
pair[] V; for(int k=0;k<6;++k) V.push(r*dir(60k));
for(int i=0;i<6;++i){
  path pet=arc(V[(i+5)%6],r,60i+60,60i+120)--arc(V[(i+1)%6],r,60i+240,60i+300)--cycle;
  fill(pet,rgb(0.7,0.9,0.7));}
for(int k=0;k<6;++k) draw(arc(V[k],r,60k+120,60k+240),darkgreen);
draw(V[0]--V[1]--V[2]--V[3]--V[4]--V[5]--cycle,linewidth(0.9));
dot((0,0)); label("$69$",(V[4]+V[5])/2,S);
\end{asy}

Find the area of the part of the lid not covered by the tea leaves.
